\documentclass[11pt]{article}
\usepackage[margin=1in]{geometry}
\begin{document}

\section*{Response to Check-In \#2 teaching-team feedback}

We appreciate the close read. In the main text we will ground the opening more in existing work on cortical processing of voice and non-voice sounds (e.g., voice-selective and speech-in-noise effects), so the motivation is not only conceptual but tied to specific studies.

The reviewers asked why \texttt{song\_only} can look stronger than \texttt{all} in the raw time-resolved decoders, which is non-intuitive if ``more electrodes'' should always help. A smaller set can be easier to fit. There are fewer features, less correlated noise, and a classifier that can focus signal without diluting it. That pattern is not, by itself, evidence that song-tagged electrodes are special. The song-selective labels come from a separate task in Norman and Haignere et al., not from the song-vs-music pair we decode here, so the analysis is not circular. Still, the fair comparison is size-matched. We address the subset-size confound with a random seven-electrode control drawn from the non-song pool (Section~5 in our pipeline), which is pre-registered in the report and notebook.

For figures and the write-up, we will define \emph{divergence} once, in the main text. It is a stimulus-level RDM, that is, a dissimilarity distance between the centroids of song and instrumental-music conditions in a sliding neural window, so it measures how separable the two classes are in representational space (with bootstrap CIs and label-shuffle nulls, as in the notebook). We will spell out what each bar and axis shows (e.g., balanced accuracy or peak divergence by subset) and avoid relying on Colab. All code and results are in the project repository. The final PDF will stand alone, with the notebook for verification only.

\end{document}
